\section{Tiered Memory Architecture}
\label{sec:memory}

Traditional browser agents operate statelessly across sessions. OpenSkynet introduces a three-tier memory system that progressively promotes knowledge from ephemeral working memory to persistent long-term storage, enabling cross-session improvement.

\subsection{Three-Tier Model}

The memory hierarchy consists of:

\begin{itemize}[leftmargin=*,itemsep=1pt]
    \item \textbf{Working Memory} (max 20 entries, 4000 characters): Holds observations and lessons from the current task. Populated by the \texttt{AutoMemory} module, which scores each observation for importance using a heuristic regex-based classifier (high-value keywords like ``always'', ``critical'', ``bug'', ``workaround'' receive +0.15; low-value keywords like ``maybe'', ``test'', ``tmp'' receive --0.10). Failed actions receive an automatic +0.2 importance boost.
    \item \textbf{Session Memory} (max 100 entries, 20000 characters): Accumulates working memory entries across a session. Entries are evicted based on a combined score: $\text{importance} \times 0.5 + \text{recency} \times 0.3 + \text{access\_frequency} \times 0.2$. High-importance entries ($\geq$ 0.7) are promoted to long-term storage.
    \item \textbf{Long-Term Memory}: Persistent dual-file storage (\texttt{MEMORY.md} for agent knowledge, \texttt{USER.md} for user preferences) using section-sign-delimited entries for human readability. Entries are written atomically via tempfile + \texttt{os.replace}. A frozen snapshot is taken at session start to preserve LLM prefix-cache invariance.
\end{itemize}

\subsection{Retrieval Architecture}

Memory retrieval employs a multi-layered approach with graceful degradation:

\begin{enumerate}[leftmargin=*,itemsep=1pt]
    \item \textbf{Vector Search:} Entries are embedded and stored in a SQLite-backed vector database with optional \texttt{sqlite-vec} extension for native ANN. The embedding provider degrades through three tiers: OpenAI \texttt{text-embedding-3-small} (1536d) $\rightarrow$ local FastEmbed \texttt{BAAI/bge-small-en-v1.5} (384d) $\rightarrow$ pure-Python TF-IDF (512d). The TF-IDF provider persists vocabulary and IDF values across sessions, enabling incremental learning.
    \item \textbf{Keyword Fallback:} When no embedding provider is available, Jaccard-like word overlap scoring over name + description + category + keywords is used.
    \item \textbf{LLM-Powered Background Review:} Every 10 turns, a background review agent analyzes recent conversation and current memory, suggesting additions, replacements, or removals. Contradiction detection uses negation-pair heuristics (e.g., ``always'' vs.\ ``never'', ``prefer'' vs.\ ``avoid'') with word-overlap similarity thresholds.
    \item \textbf{Context Filtering:} Task-specific context injection uses a multi-factor ranking: text relevance (0.5) + recency (0.25) + access frequency (0.15) + entry-type bonus (0.10--0.15), ensuring only the most pertinent memories consume context window budget.
\end{enumerate}

\subsection{Consolidation}

When memory limits are reached, the \texttt{MemoryConsolidator} applies three strategies in order: (1) merge small entries ($<$ 60 characters) with semicolon separators; (2) summarize old entries by truncating to the first half of sentences; and (3) remove the lowest-scored entries. An optional LLM-based consolidation path summarizes the oldest entries when they account for at least 50\% of the needed space, avoiding wasted LLM calls on trivial memory pressure.

\subsection{Security}

All memory writes pass through a content security scanner that detects prompt injection attempts, credential leaks (\texttt{api\_key}, \texttt{token}, \texttt{secret}), backdoor commands (\texttt{authorized\_keys}), destructive operations (\texttt{rm -rf /}, \texttt{DROP TABLE}), and invisible Unicode characters (zero-width, bidi overrides, format characters). A streaming context scrubber removes \texttt{<memory-context>} tags from LLM output to prevent internal memory state from leaking to users.
